\chapter{Reference problem, conventions, and the master formula}
\label{chap:reference}

\section{Scope}

The unperturbed electric fields and the unperturbed specular reflection
amplitude are assumed to be available from the layered-field solver.  The
problem addressed here is the remaining atomistic matrix element for finite
unit cells in the graded interface region.  The terminal semi-infinite bulk is
included only to show that the new formula reduces to the already implemented
bulk result.

The derivation is electromagnetic throughout.  Following Renaud, Lazzari, and
Leroy, the incident field and the reciprocal field launched from the detector
are denoted by electric fields $\bm E_0$; no quantum-mechanical wavefunction
notation is used.  The medium is nonmagnetic, isotropic, linear, and reciprocal.
Magnetic and tensor-resonant scattering are outside the present scope.

This chapter derives, rather than quotes, both the matrix element and the
factor that converts it into a measurable rod amplitude.  The chain is: Born
expansion, Green function of the stratified reference built from two
homogeneous solutions and their Wronskian, lattice expansion of the index
contrast, and lateral collapse.  No far-field (Fraunhofer) approximation over
the illuminated footprint is used at any point; Sec.~\ref{sec:why-not-farfield}
explains why one must not be.

\section{Fourier, surface, and refractive-index conventions}

We use the crystallographic space-time convention
\begin{equation}
  \bm E(\bm r,t)\propto
  \exp\!\left[\ii\left(\omega t-\bm k\mathbin{\cdot}\bm r\right)\right],
  \qquad
  \kzero=\frac{2\pi}{\lambda},
  \qquad
  \bm Q=\bm k_f-\bm k_i,
  \label{eq:convention}
\end{equation}
and the Fourier transform
\begin{equation}
  \widetilde g(\bm Q)
  =\int g(\bm r)\exp(\ii\bm Q\mathbin{\cdot}\bm r)\,\dd^3r.
  \label{eq:fourier}
\end{equation}
The $z$ axis is the outward surface normal: vacuum is at large positive $z$
and the bulk extends toward negative $z$.  Momentum transfer is in
Angstrom$^{-1}$ and all real-space coordinates below are in Angstrom.

The prepared optical reference consists of planar media $m=0,\ldots,M$.  In
medium $m$,
\begin{equation}
  n_m=1-\delta_m-\ii\beta_m,
  \qquad \delta_m,\beta_m\in\mathbb R,
  \label{eq:index}
\end{equation}
and the medium occupies the interval
\begin{equation}
  I_m=(z_m^-,z_m^+).
\end{equation}
One or both endpoints may be infinite for a terminal medium.  A fixed
half-open convention must be used in code when an atom lies exactly on a
boundary.

orGUI uses the negative-root Renaud normal wavevector
\begin{equation}
  k_{z,m}=-\kzero
  \sqrt{n_m^2-n_0^2\cos^2\alpha},
  \qquad
  \operatorname{Re}k_{z,m}\leq0,
  \quad
  \operatorname{Im}k_{z,m}\geq0.
  \label{eq:renaud-kz}
\end{equation}
The incident and reciprocal exit solutions are calculated in the same frozen
array of complex refractive indices and at the same physical boundaries.

Equation~\eqref{eq:renaud-kz} fixes what ``upward'' means, and every sign in
this document follows from it.  Because $\operatorname{Re}k_{z,m}\leq0$, a
wave travelling towards $+z$ --- upward, out of the sample and towards the
detector --- carries $\exp(+\ii k_{z,m}z)$, while a wave travelling towards
$-z$, downward onto the sample, carries $\exp(-\ii k_{z,m}z)$.  This is the
opposite pairing from the one a positive-root convention would give, and it is
the single most common source of sign errors when comparing with the
literature.  The whole of Secs.~\ref{sec:born}--\ref{sec:collapse} is carried
out in Eq.~\eqref{eq:renaud-kz}'s convention, so no conversion is ever needed.

\section{Scattering-factor density and the index contrast}

For atom $a$, orGUI uses
\begin{equation}
  f_a(Q^2,E)=f^0_{a,\mathrm{WK}}(Q^2)+f'_a(E)+\ii f''_a(E),
  \label{eq:atomic-factor}
\end{equation}
where the Waasmaier--Kirfel representation is evaluated analytically at the
generally complex channel value $Q^2$.  Its zero-transfer value, including the
same anomalous terms, is also used to generate the optical profile.  Replacing
that value by an independent atomic-number approximation would break the
reference-density cancellation.  Equation~\eqref{eq:atomic-factor} is the only
quantity in this document called $f$.

Let $\rho_f(\bm r)$ be the complex scattering-factor density in electrons per
Angstrom$^3$.  To first order in the X-ray susceptibility,
\begin{equation}
  n^2(\bm r)\simeq
  1-\frac{4\pi r_e}{\kzero^2}\rho_f(\bm r),
  \label{eq:index-density}
\end{equation}
with $r_e$ the classical electron radius.  The piecewise-constant reference
density represented by Eq.~\eqref{eq:index} is therefore
\begin{equation}
  \boxed{
  \overline\rho_{f,m}
  =\frac{\kzero^2}{2\pi r_e}
   \left(\delta_m+\ii\beta_m\right).
  }
  \label{eq:reference-density}
\end{equation}
Write $n_0^2(z)=\sum_m n_m^2\,\mathbf 1_{I_m}(z)$ for the prepared, laterally
uniform reference and
\begin{equation}
  \delta n^2(\bm r)=n^2(\bm r)-n_0^2(z)
  =-\frac{4\pi r_e}{\kzero^2}\,\Delta\rho_f(\bm r),
  \qquad
  \Delta\rho_f(\bm r)
  =\rho_f(\bm r)-\sum_m\overline\rho_{f,m}\,\mathbf 1_{I_m}(z),
  \label{eq:density-contrast}
\end{equation}
for the index contrast that drives all scattering.  Equations~\eqref{eq:index}
and \eqref{eq:reference-density} enter the problem only through $n_0^2(z)$, and
Sec.~\ref{sec:collapse} shows exactly where they survive.

\section{The Born expansion}
\label{sec:born}

At the small scattering angles of interest the vectorial propagation equation
reduces to a scalar Helmholtz equation for the component of the field along
the analysed polarization (Renaud \S5.3.6).  Writing $E$ for that component,
\begin{equation}
  \left[\nabla^2+\kzero^2n^2(\bm r)\right]E(\bm r)=0,
  \qquad
  n^2(\bm r)=n_0^2(z)+\delta n^2(\bm r),
  \label{eq:helmholtz}
\end{equation}
so that
\begin{equation}
  \left[\nabla^2+\kzero^2n_0^2(z)\right]E(\bm r)
  =-\kzero^2\,\delta n^2(\bm r)\,E(\bm r).
  \label{eq:driven}
\end{equation}
Defining the Green function of the \emph{reference} problem by
\begin{equation}
  \left[\nabla^2+\kzero^2n_0^2(z)\right]G(\bm r,\bm r')=\delta(\bm r-\bm r'),
  \label{eq:green-def}
\end{equation}
Eq.~\eqref{eq:driven} is equivalent to the integral equation
\begin{equation}
  E(\bm r)=\bm E_0(\bm r)
  -\kzero^2\!\int\!\dd^3r'\,G(\bm r,\bm r')\,\delta n^2(\bm r')\,E(\bm r'),
  \label{eq:lippmann}
\end{equation}
whose iteration is the Born series.  Retaining the first order, with $\bm E_0$
and $G$ both those of the \emph{stratified} reference, gives the distorted-wave
Born approximation:
\begin{equation}
  \boxed{
  \delta E(\bm r)
  =-\kzero^2\!\int\!\dd^3r'\,G(\bm r,\bm r')\,\delta n^2(\bm r')\,
   \bm E_0(\bm r',\bm k_i).
  }
  \label{eq:dwba-master}
\end{equation}
Reflection, transmission, refraction, and repeated reflection at the planar
interfaces are already contained in $\bm E_0$ and in $G$; the perturbation is
scattered exactly once.

\paragraph{Sign remark.}  Renaud's Eqs.~(71) and (73) as printed carry
$+\kzero^2\delta n^2$ on the right-hand side.  With his Eq.~(70)'s definition
$\delta n^2=n^2-n_0^2$, adopted here in Eq.~\eqref{eq:density-contrast}, the
sign is as in Eqs.~\eqref{eq:driven} and \eqref{eq:lippmann}.  The overall sign
of the final result is verified independently against the Born--Fresnel limit
in Sec.~\ref{sec:born-check}.

\section{The Green function of the stratified reference}
\label{sec:green}

This section does the work that is usually skipped.  Nothing in it is
approximate.

\subsection{Separation into a one-dimensional problem}

Because $n_0^2$ depends only on $z$, Eq.~\eqref{eq:green-def} separates under a
lateral Fourier transform:
\begin{equation}
  G(\bm r,\bm r')
  =\frac{1}{(2\pi)^2}\int\dd^2q\;g_q(z,z')\,
   \ee^{-\ii\bm q\cdot(\rpar-\rpar')},
  \qquad
  \left[\frac{\dd^2}{\dd z^2}+\kzero^2n_0^2(z)-q^2\right]g_q(z,z')
  =\delta(z-z').
  \label{eq:separation}
\end{equation}
The lateral wavevector $\bm q$ enters the one-dimensional operator only through
the constant $-q^2$; the normal wavevector belonging to $\bm q$ in medium $m$
is Eq.~\eqref{eq:renaud-kz}'s $k_{z,m}$.

\subsection{Wronskian construction}

Equation~\eqref{eq:separation} is a one-dimensional Green-function problem for
an operator with no first-derivative term.  Let $u_-$ and $u_+$ solve the
homogeneous equation subject to one physical boundary condition each,
\begin{align}
  u_-(z):&\quad\text{no wave incoming from below, as }z\to-\infty,
  \label{eq:bc-minus}\\
  u_+(z):&\quad\text{no wave incoming from above, as }z\to+\infty,
  \label{eq:bc-plus}
\end{align}
each fixed up to one multiplicative constant.  Then
\begin{equation}
  g_q(z,z')=\frac{u_-(z_<)\,u_+(z_>)}{W},
  \qquad
  W=u_-u_+'-u_-'u_+,
  \label{eq:wronskian-construction}
\end{equation}
with $z_<=\min(z,z')$ and $z_>=\max(z,z')$.  The Wronskian is independent of
$z$,
\begin{equation}
  \frac{\dd W}{\dd z}
  =u_-u_+''-u_-''u_+
  =u_-\!\left(q^2-\kzero^2n_0^2\right)u_+
   -\left(q^2-\kzero^2n_0^2\right)u_-u_+=0,
  \label{eq:wronskian-constant}
\end{equation}
precisely because the operator has no $\dd/\dd z$ term, so it may be evaluated
wherever it is easiest.  The normalisation of $u_\pm$ cancels in
Eq.~\eqref{eq:wronskian-construction}.

\subsection{What the two solutions are physically}

Using the convention established after Eq.~\eqref{eq:renaud-kz}, in vacuum
($m=0$):
\begin{itemize}
  \item $u_+$ is fixed by Eq.~\eqref{eq:bc-plus}: above the whole structure it
        is purely upward-going,
        \begin{equation}
          u_+(z)=\ee^{+\ii k_{z,0}z}.
          \label{eq:u-plus}
        \end{equation}
        How it continues downward through the stack is never needed, because
        the observation point is always in vacuum.
  \item $u_-$ is fixed by Eq.~\eqref{eq:bc-minus}: deep in the terminal medium
        it carries only the transmitted, downward or evanescently decaying
        wave, which is the physical condition $A_M^-=0$.  That single condition,
        continued \emph{upward} through the stack, produces in vacuum a
        superposition of a downward and an upward component; normalising the
        downward one to unity,
        \begin{equation}
          u_-(z)=\ee^{-\ii k_{z,0}z}+r_0\,\ee^{+\ii k_{z,0}z}
          \qquad\text{(in vacuum)}.
          \label{eq:u-minus}
        \end{equation}
        This is the ordinary Fresnel/Parratt reference field for a plane wave
        incident from above --- incident plus reflected, with only a
        transmitted wave below --- and $r_0$ is the unperturbed specular
        reflection amplitude of Eq.~\eqref{eq:specular-reflectivity}.  The
        condition that \emph{defines} $u_-$ is imposed at $z\to-\infty$, but
        the solution it selects is the familiar one driven from above.
\end{itemize}
Since the exit wavevector $\bm k_f$ points away from the sample, $-\bm k_f$
points onto it, and $u_-$ evaluated for that direction is exactly the
reciprocal exit field,
\begin{equation}
  u_-(z')=\bm E_0(z',-\bm k_f).
  \label{eq:u-minus-identification}
\end{equation}
This is the content of the reciprocity theorem, obtained here by solving a
boundary-value problem rather than by propagating a field to a distant
detector.  It is why the reciprocal exit field is never complex conjugated:
it is a solution of the same homogeneous equation as the incident field, not
its adjoint.

\subsection{The Wronskian, and the Green function}

Evaluate $W$ in vacuum from Eqs.~\eqref{eq:u-plus}--\eqref{eq:u-minus}:
\begin{align}
  W&=\left[\ee^{-\ii k_{z,0}z}+r_0\ee^{\ii k_{z,0}z}\right]
     \left(\ii k_{z,0}\ee^{\ii k_{z,0}z}\right)
    -\left[-\ii k_{z,0}\ee^{-\ii k_{z,0}z}
       +\ii k_{z,0}r_0\ee^{\ii k_{z,0}z}\right]\ee^{\ii k_{z,0}z}\nonumber\\
   &=\ii k_{z,0}+\ii k_{z,0}r_0\ee^{2\ii k_{z,0}z}
     +\ii k_{z,0}-\ii k_{z,0}r_0\ee^{2\ii k_{z,0}z}
   \;=\;2\ii k_{z,0}.
  \label{eq:wronskian-value}
\end{align}
The $r_0$ terms cancel identically: \emph{the Wronskian does not depend on the
stack}, only on the vacuum $k_{z,0}$.  This one fact is the origin of every
factor $1/k_z$ in this subject, including the one that appears without
explanation in the reflectivity literature.

With the observation point above all sources, $z>z'$, so $z_<=z'$, $z_>=z$, and
Eqs.~\eqref{eq:wronskian-construction}, \eqref{eq:u-plus},
\eqref{eq:u-minus-identification} and \eqref{eq:wronskian-value} give
\begin{equation}
  \boxed{
  g_q(z,z')
  =\frac{\bm E_0(z',-\bm k_f)\,\ee^{\ii k_{z,0}z}}{2\ii k_{z,0}}
  =-\frac{\ii}{2k_{z,0}}\,\ee^{\ii k_{z,0}z}\,\bm E_0(z',-\bm k_f).
  }
  \label{eq:green-1d}
\end{equation}
For a homogeneous reference, $r_0=0$ and $u_-(z')=\ee^{-\ii k_{z,0}z'}$, so
Eq.~\eqref{eq:green-1d} collapses to
$g_q=-(\ii/2k_{z,0})\ee^{\ii k_{z,0}|z-z'|}$, whose derivative jump across
$z=z'$ is unity; substituting that into Eq.~\eqref{eq:separation} reproduces
the Weyl angular-spectrum representation of an outgoing spherical wave.  The
Weyl expansion is therefore not an extra ingredient but the vacuum special case
of Eq.~\eqref{eq:green-1d}.

\subsection{Why the far-field route is not available here}
\label{sec:why-not-farfield}

The textbook shortcut replaces $G$ by its far-field limit, expanding
$|\bm r-\bm r'|\simeq r-\bm k_f\cdot\bm r'/\kzero$ and pulling a single
spherical wave $\ee^{-\ii\kzero r}/r$ out of the integral.  That expansion
requires the detector distance $R$ to exceed the Fraunhofer distance of the
\emph{whole illuminated footprint}.  For a grazing-incidence footprint of
lateral size $L\sim1$~mm to 1~cm at $\lambda\sim1$~\AA, the neglected
quadratic phase is
\begin{equation}
  \delta(\text{phase})\sim\kzero\frac{L^2}{R}\sim10^{4}\text{--}10^{7}\gg1,
  \label{eq:fraunhofer-failure}
\end{equation}
so the far-field form of $G$ is never valid on the specular rod.
Equation~\eqref{eq:green-1d} avoids the issue entirely: it is exact at any $R$.

The polarization factor is a separate matter.  The same far-field substitution
is conventionally used to fix the exit polarization vector, but there it enters
undifferentiated and is \emph{not} multiplied by $\kzero$, so its error is only
$\sim L/R\sim10^{-2}$ rather than Eq.~\eqref{eq:fraunhofer-failure}.  Fixing the
polarization basis at the nominal $\bm k_f$, as
Eq.~\eqref{eq:polarization-contractions} does, therefore remains accurate even
where the phase treatment above is mandatory.

\section{Lattice expansion and the lateral collapse}
\label{sec:collapse}

\subsection{The index contrast along the surface lattice}

Parallel to the surface, $n^2(\bm r)$ is periodic on the two-dimensional
lattice of the reference lateral cell, of area $A_{\rm ref}$, while $n_0^2(z)$
is laterally uniform.  Their difference $\delta n^2$ is therefore
two-dimensional-periodic and has the Fourier series
\begin{equation}
  \delta n^2(\rpar,z)=\sum_{\bm h}\delta n^2_{\bm h}(z)\,
  \ee^{-\ii\bm h\cdot\rpar},
  \qquad
  \delta n^2_{\bm h}(z)
  =\frac{1}{A_{\rm ref}}\int_{\rm cell}\dd^2r_\parallel\;
   \delta n^2(\rpar,z)\,\ee^{\ii\bm h\cdot\rpar},
  \label{eq:lattice-expansion}
\end{equation}
where $\bm h$ runs over the reciprocal lattice of that cell, i.e.\ over the
crystal truncation rods $(h,k)$.  Because $n_0^2(z)$ is laterally uniform it
contributes to Eq.~\eqref{eq:lattice-expansion} at $\bm h=\bm 0$ and nowhere
else:
\begin{align}
  \delta n^2_{\bm h}(z)
  &=\frac{1}{A_{\rm ref}}\int_{\rm cell}\dd^2r_\parallel\;n^2(\rpar,z)\,
    \ee^{\ii\bm h\cdot\rpar},
  &&\bm h\neq\bm0,
  \label{eq:coeff-nonspec}\\
  \delta n^2_{\bm 0}(z)
  &=\frac{1}{A_{\rm ref}}\int_{\rm cell}\dd^2r_\parallel\;n^2(\rpar,z)
    \;-\;n_0^2(z),
  &&\bm h=\bm 0.
  \label{eq:coeff-spec}
\end{align}
Equations~\eqref{eq:coeff-nonspec}--\eqref{eq:coeff-spec} are the entire
content of the specular/non-specular distinction, and they are the reason this
document needs only one matrix element.  The $\delta_m,\beta_m$ of the prepared
profile are not something subtracted afterwards from an atomic sum: they are
part of the $\bm h=\bm 0$ Fourier coefficient, and part of nothing else.

\subsection{Collapse}

Insert Eqs.~\eqref{eq:separation}, \eqref{eq:green-1d} and
\eqref{eq:lattice-expansion} into Eq.~\eqref{eq:dwba-master}.  Stratification
lets the lateral phase of the incident field be factored out,
$\bm E_0(\bm r',\bm k_i)
=\ee^{-\ii\bm k_{i\parallel}\cdot\rpar'}\bm E_0(z',\bm k_i)$, so all
$\rpar'$ dependence sits in three exponentials and the lateral integral is
elementary:
\begin{equation}
  \int\dd^2r_\parallel'\;
  \ee^{\ii\bm q\cdot\rpar'}\,
  \ee^{-\ii\bm h\cdot\rpar'}\,
  \ee^{-\ii\bm k_{i\parallel}\cdot\rpar'}
  =(2\pi)^2\,\delta^{(2)}\!\left(\bm q-\bm k_{i\parallel}-\bm h\right).
  \label{eq:lateral-delta}
\end{equation}
The delta function is the Laue condition: for each rod exactly one
angular-spectrum mode survives,
\begin{equation}
  \bm q=\bm k_{i\parallel}+\bm h\equiv\bm k_{f\parallel},
  \qquad\text{that is}\qquad
  \Qpar=\bm h,
  \label{eq:laue}
\end{equation}
at any incidence angle, any exit angle, and any azimuth.  The $1/r$ spherical
falloff has cancelled against the density of angular-spectrum modes: what
emerges is a genuine, non-decaying plane wave per rod.

\subsection{The matrix element and the rod amplitude}

Collecting the constants of Eqs.~\eqref{eq:dwba-master}, \eqref{eq:separation},
\eqref{eq:green-1d} and \eqref{eq:lateral-delta},
\begin{equation}
  -\kzero^2\cdot\frac{1}{(2\pi)^2}
  \cdot\left(-\frac{\ii}{2k_{z,0,f}}\right)\cdot(2\pi)^2
  =+\frac{\ii\kzero^2}{2k_{z,0,f}},
  \label{eq:constants}
\end{equation}
and converting from $\delta n^2$ to $\Delta\rho_f$ with
Eq.~\eqref{eq:density-contrast}, the $\kzero^2$ cancels exactly.  Defining the
matrix element of rod $\bm h$, normalised to one reference lateral cell,
\begin{equation}
  \boxed{
  F_{\bm h}
  =\int\dd z\int_{\rm cell}\dd^2r_\parallel\;
   \Delta\rho_f(\rpar,z)\,
   \ee^{\ii\bm h\cdot\rpar}\,
   \bm E_0(z,-\bm k_f)
   \mathbin{\cdot}
   \bm E_0(z,\bm k_i),
  }
  \label{eq:Fh-definition}
\end{equation}
the scattered field on rod $\bm h$ is the plane wave
\begin{equation}
  \boxed{
  r_{\bm h}
  =-\frac{2\pi\ii r_e}{A_{\rm ref}\,k_{z,0,f}}\,F_{\bm h}.
  }
  \label{eq:rod-amplitude}
\end{equation}

Equation~\eqref{eq:Fh-definition} is the only matrix element in this document.
It is used unchanged for every rod:
\begin{itemize}
  \item for $\bm h\neq\bm 0$ its integrand contains no reference term, by
        Eq.~\eqref{eq:coeff-nonspec};
  \item for $\bm h=\bm 0$ its integrand contains the constant
        $\delta_m+\ii\beta_m$ density of every prepared slab, by
        Eq.~\eqref{eq:coeff-spec}, through $\Delta\rho_f$'s
        $-\sum_m\overline\rho_{f,m}\mathbf 1_{I_m}$ term.
\end{itemize}
Chapter~\ref{chap:unit-cells} evaluates Eq.~\eqref{eq:Fh-definition} for atoms,
producing the four coherent channels.  Chapter~\ref{chap:specular} evaluates
the reference part of $F_{\bm 0}$ and partitions it among generated records.
The reciprocal exit field in Eq.~\eqref{eq:Fh-definition} is not complex
conjugated; conjugation enters only after the coherent amplitude has been
formed.

\subsection{Born--Fresnel check}
\label{sec:born-check}

Take a uniform half-space $z<0$ of density $\rho_f$ against a vacuum
reference, so that only $\bm h=\bm 0$ contributes, the field product reduces
to $\ee^{\ii Q_zz}$, and
$\int_{-\infty}^{0}\ee^{\ii Q_zz}\dd z=1/(\ii Q_z)$.  Then
$F_{\bm 0}=A_{\rm ref}\rho_f/(\ii Q_z)$ and, with $Q_z=-2k_{z,0}>0$ from
Eq.~\eqref{eq:channel-qz},
\begin{equation}
  r_{\bm 0}
  =-\frac{2\pi\ii r_e}{A_{\rm ref}k_{z,0}}\cdot
    \frac{A_{\rm ref}\rho_f}{\ii Q_z}
  =-\frac{2\pi r_e\rho_f}{k_{z,0}Q_z}
  =+\frac{4\pi r_e\rho_f}{Q_z^{2}}\;>0,
  \label{eq:born-fresnel}
\end{equation}
which is the known kinematic reflectivity amplitude, with the positive sign of
the exact Fresnel expansion $r=(k_{z,0}-k_{z,1})/(k_{z,0}+k_{z,1})$ for $n<1$.
Both the sign and the constant in Eq.~\eqref{eq:rod-amplitude} are therefore
fixed independently of any external sign convention.
